\chapter{La integral de Lebesgue}\label{ch:b3:lebesgue}

La integral de Riemann corta el \emph{dominio} en intervalos pequeños;
la de Lebesgue corta el \emph{recorrido}: para integrar $f$, se
mide el conjunto $\{f > t\}$. El cambio
parece inocente y es revolucionario. Los límites y las integrales,
eternamente enfrentados en la teoría de Riemann (¡se exigía convergencia
uniforme!), quedan reconciliados por tres teoremas de convergencia ---
convergencia monótona, Fatou y convergencia dominada --- cuyas hipótesis
son casi vergonzosamente débiles. Este capítulo construye la integral
sobre un \omterm{def:b3:measure:measure}{espacio de medida} arbitrario
$(X, \mathcal A,
\mu)$, demuestra los tres teoremas, zanja la relación exacta con la
integral de Riemann (una función acotada es
\omterm{def:b3:lebesgue:l1}{integrable} Riemann si y solo si es
\omterm{def:b3:topology:continuity}{continua}
\omterm{def:b3:lebesgue:l1}{en casi todo punto}) e industrializa la
derivación de integrales dependientes de un parámetro --- la técnica que
el problema de fin de semana usa para calcular
$\int_0^\infty\frac{\sin x}x\,\dd x$ y
$\int_\R \eu^{-x^2}\dd x$.

\section{Funciones medibles}

\begin{definition}\label{def:b3:lebesgue:measurable}
Sean $(X, \mathcal A)$ y $(Y, \mathcal B)$ espacios medibles.
$f \colon X \to Y$ es \emph{medible}\index{función medible} si
$f^{-1}(B) \in \mathcal A$ para todo $B \in
\mathcal B$. Para funciones reales (o con valores en
$[-\infty,+\infty]$), $Y = \R$ lleva su $\sigma$-álgebra de Borel, y
basta comprobar $f^{-1}(\intoo t{+\infty}) = \{f > t\} \in
\mathcal A$ para todo $t \in \R$: los buenos conjuntos $\{B :
f^{-1}(B) \in \mathcal A\}$ forman una $\sigma$-álgebra (las preimágenes
conmutan con las operaciones conjuntistas) que contiene las semirrectas
generadoras (la~\cref{def:b3:measure:borel},
\cref{met:b3:measure:goodsets}).
\end{definition}

\begin{proposition}\label{prop:b3:lebesgue:stability}
(a) Las composiciones de aplicaciones
\omterm{def:b3:lebesgue:measurable}{medibles} son
\omterm{def:b3:lebesgue:measurable}{medibles}; las
aplicaciones \omterm{def:b3:topology:continuity}{continuas} son
\omterm{def:b3:lebesgue:measurable}{medibles} Borel. (b) Si
$f, g \colon X \to \R$ son
\omterm{def:b3:lebesgue:measurable}{medibles}, también lo son $f + g$,
$fg$, $\max(f,g)$, $\abs f$, $\lambda f$.
(c) Si $(f_n)$ son \omterm{def:b3:lebesgue:measurable}{medibles} con
valores en $[-\infty, +\infty]$, entonces $\sup_nf_n$, $\inf_nf_n$,
$\limsup f_n$, $\liminf f_n$ son
\omterm{def:b3:lebesgue:measurable}{medibles}; y si $f_n \to f$
puntualmente, $f$ es \omterm{def:b3:lebesgue:measurable}{medible}.
\end{proposition}

\begin{proof}
(a) $(g\circ f)^{-1}(B) = f^{-1}(g^{-1}(B))$; la
\omterm{def:b3:topology:continuity}{continuidad} da la
\omterm{def:b3:lebesgue:measurable}{medibilidad} a través de los
abiertos generadores
(el~\cref{pb:b3:measure:1}, pregunta 10, en forma general). (b)
$(f, g) \colon X \to \R^2$ es
\omterm{def:b3:lebesgue:measurable}{medible} para la $\sigma$-álgebra de
Borel de $\R^2$ --- compruébese sobre las cajas abiertas, que generan
(los abiertos de $\R^2$ son uniones numerables de cajas racionales):
$(f,g)^{-1}(U\times V) = f^{-1}(U)\cap g^{-1}(V)$ --- y
$+, \times, \max$ son \omterm{def:b3:topology:continuity}{continuas}
$\R^2 \to \R$: compóngase.
(c) $\{\sup f_n > t\} = \bigcup_n\{f_n > t\}$; $\inf = -\sup(-)$;
$\limsup = \inf_N\sup_{n \geq N}$; y un límite puntual es su propio
$\limsup$.
\end{proof}

\begin{definition}\label{def:b3:lebesgue:simple}
Una \emph{función simple}\index{función simple} es una
\omterm{def:b3:lebesgue:measurable}{función medible} con un número
finito de valores: $s = \sum_{i=1}^n
c_i\,\mathbf 1_{A_i}$, con $A_i \in \mathcal A$ disjuntos y $c_i \geq
0$ (para la teoría no negativa). Su integral es
\[
\int s\,\dd\mu = \sum_i c_i\,\mu(A_i) \in [0, +\infty]
\]
(con el convenio $0\cdot\infty = 0$); el valor no depende de la
\omterm{def:b3:representations:rep}{representación} (refínense dos
particiones).
\end{definition}

\begin{theorem}[Aproximación por funciones simples]
\label{thm:b3:lebesgue:approximation}
Toda \omterm{def:b3:lebesgue:measurable}{función medible}
$f \colon X \to [0, +\infty]$ es límite puntual de una sucesión
\emph{creciente} de \omterm{def:b3:lebesgue:simple}{funciones simples}:
\[
s_n = \sum_{k=1}^{n2^n} \frac{k-1}{2^n}\,
\mathbf 1_{\{\frac{k-1}{2^n} \leq f < \frac k{2^n}\}}
+ n\,\mathbf 1_{\{f \geq n\}} \nearrow f .
\]
\end{theorem}

\begin{proof}
Cada $s_n$ es simple (los conjuntos son preimágenes de conjuntos de
Borel). Monotonía: al pasar de $n$ a $n+1$, cada nivel diádico se parte
en dos y el valor asignado nunca decrece (un punto con
$\frac{k-1}{2^n} \leq f(x) < \frac k{2^n}$ recibe $\frac{2k-2}{2^{n+1}}$
o $\frac{2k-1}{2^{n+1}}$, ambos $\geq
\frac{k-1}{2^n}$; y el tope $n$ también sube). Convergencia: si
$f(x) < \infty$, para $n > f(x)$ se tiene $f(x) - s_n(x) \leq
2^{-n}$; si $f(x) = \infty$, $s_n(x) = n \to \infty$.
\end{proof}

\section{La integral y los teoremas de convergencia}

\begin{definition}\label{def:b3:lebesgue:integral}
Para $f \geq 0$ \omterm{def:b3:lebesgue:measurable}{medible}:
\[
\int f \,\dd\mu = \sup\Bigl\{\int s\,\dd\mu : s \text{ simple},
\ 0 \leq s \leq f\Bigr\} \in [0, +\infty].
\]
Es monótona en $f$ por construcción y extiende el caso simple (para $f$
simple, el supremo se alcanza en $f$: compárense integrales de \omterm{def:b3:lebesgue:simple}{funciones
simples} mediante refinamientos comunes).
\end{definition}

\begin{theorem}[Convergencia monótona, Beppo Levi]
\label{thm:b3:lebesgue:mct}
Si $0 \leq f_n \nearrow f$ puntualmente (con las
\omterm{def:b3:lebesgue:measurable}{medibles}), entonces
\[
\int f_n\,\dd\mu \nearrow \int f\,\dd\mu .
\]
\index{teorema de la convergencia monótona}
\end{theorem}

\begin{proof}
$f$ es \omterm{def:b3:lebesgue:measurable}{medible}
(la~\cref{prop:b3:lebesgue:stability}(c)) y $\int f_n$ crece hacia
cierto $L \leq \int f$ (monotonía). Recíprocamente, fíjense una
$s = \sum c_i\mathbf 1_{A_i} \leq f$ simple y $\theta \in (0,1)$; los
conjuntos $E_n = \{f_n \geq \theta s\}$ son
\omterm{def:b3:lebesgue:measurable}{medibles} y crecen hacia $X$ (donde
$s(x) > 0$: $f(x)
\geq s(x) > \theta s(x)$, luego a la larga $f_n(x) \geq \theta
s(x)$; donde $s(x) = 0$: trivialmente). Entonces
\[
\int f_n \geq \int_{E_n}\theta s\,\dd\mu
= \theta\sum_i c_i\,\mu(A_i \cap E_n)
\xrightarrow[n\to\infty]{} \theta\sum_ic_i\,\mu(A_i)
= \theta\int s
\]
por \omterm{def:b3:topology:continuity}{continuidad} por abajo
(la~\cref{prop:b3:measure:basics}(c)). Así, $L \geq \theta\int s$ para
todo $\theta < 1$ y toda $s
\leq f$ simple: $L \geq \int f$.
\end{proof}

\begin{corollary}\label{cor:b3:lebesgue:additivity}
Para $f, g \geq 0$ \omterm{def:b3:lebesgue:measurable}{medibles} y
$c \geq 0$: $\int(f + g) =
\int f + \int g$ y $\int cf = c\int f$; y para una serie de
funciones \omterm{def:b3:lebesgue:measurable}{medibles} no negativas,
$\int\sum_nf_n =
\sum_n\int f_n$.
\end{corollary}

\begin{proof}
Para las \omterm{def:b3:lebesgue:simple}{funciones simples}, la
aditividad es un cálculo sobre un refinamiento común. En general,
tómense $s_n \nearrow f$, $t_n \nearrow g$
(el~\cref{thm:b3:lebesgue:approximation}): $s_n + t_n \nearrow f +
g$, y el teorema de la convergencia monótona pasa la aditividad al
límite. El enunciado sobre series es ese teorema aplicado a las sumas
parciales.
\end{proof}

\begin{theorem}[Lema de Fatou]\label{thm:b3:lebesgue:fatou}
Para $f_n \geq 0$ \omterm{def:b3:lebesgue:measurable}{medibles}:
\[
\int \liminf_n f_n \,\dd\mu \;\leq\; \liminf_n \int
f_n\,\dd\mu .
\]
\index{lema de Fatou}
\end{theorem}

\begin{proof}
Sea $g_N = \inf_{n\geq N}f_n$: es
\omterm{def:b3:lebesgue:measurable}{medible}, $0 \leq g_N \nearrow
\liminf f_n$, y $g_N \leq f_n$ para todo $n \geq N$, luego
$\int g_N \leq \inf_{n \geq N}\int f_n$. Aplíquese la convergencia
monótona al miembro izquierdo: $\int\liminf f_n = \lim_N\int g_N \leq
\lim_N\inf_{n\geq N}\int f_n = \liminf\int f_n$.
\end{proof}

\begin{definition}\label{def:b3:lebesgue:l1}
Una $f \colon X \to \R$ \omterm{def:b3:lebesgue:measurable}{medible} (o
$\C$) es \emph{integrable}\index{función integrable} si $\int\abs
f\,\dd\mu < \infty$; entonces $\int f = \int f^+ - \int f^-$ (partes
positiva y negativa; partes real e imaginaria en el caso complejo). La
integral es lineal sobre las funciones integrables (descompóngase y
recombínense partes positivas; el caso complejo se reduce al real) y
cumple $\abs{\int f} \leq
\int\abs f$ (caso real: $\pm\int f = \int(\pm f) \leq
\int\abs f$; caso complejo: multiplíquese por una constante unimodular
para hacer real la integral). Una propiedad se cumple \emph{en casi todo
punto} (c.t.p.)\index{en casi todo punto} si solo falla sobre un
conjunto $\mu$-nulo; modificar $f$ sobre un conjunto nulo no cambia
ninguna integral (la diferencia está dominada por $\infty\cdot\mathbf
1_N$, de integral $0$).
\end{definition}

\begin{theorem}[Convergencia dominada]\label{thm:b3:lebesgue:dct}
Sea $f_n \to f$ \omterm{def:b3:lebesgue:l1}{en casi todo punto}, con $\abs{f_n} \leq g$ \omterm{def:b3:lebesgue:l1}{en casi todo
punto} para una $g$ \emph{\omterm{def:b3:lebesgue:l1}{integrable}} fija.
Entonces $f$ es \omterm{def:b3:lebesgue:l1}{integrable} y
\[
\int f_n\,\dd\mu \longrightarrow \int f\,\dd\mu,
\qquad\text{de hecho}\quad \int\abs{f_n - f}\,\dd\mu \to 0 .
\]
\index{teorema de la convergencia dominada}
\end{theorem}

\begin{proof}
Descártese un conjunto nulo para que las hipótesis sean puntuales.
$\abs f
\leq g$: $f$ es \omterm{def:b3:lebesgue:l1}{integrable}. Las funciones
$h_n = 2g - \abs{f_n
- f} \geq 0$ cumplen $\liminf h_n = 2g$; Fatou da
\[
\int 2g \leq \liminf\int\bigl(2g - \abs{f_n - f}\bigr)
= \int 2g - \limsup\int\abs{f_n - f},
\]
luego $\limsup\int\abs{f_n - f} \leq 0$ (la resta es lícita:
$\int 2g < \infty$). Por último, $\abs{\int f_n - \int f} \leq
\int\abs{f_n - f} \to 0$.
\end{proof}

\begin{method}\label{met:b3:lebesgue:convergence}
Ante $\lim_n\int f_n$, pruébese, por este orden: (1) ¿es la sucesión
monótona (o una serie de términos no negativos)? Convergencia monótona,
sin necesidad de \omterm{def:b3:lebesgue:l1}{integrabilidad}. (2) ¿Hay
un único dominador \omterm{def:b3:lebesgue:l1}{integrable}
$g \geq \abs{f_n}$, hallado con cotas groseras («$\sup_n$» las
estimaciones)? Convergencia dominada. (3) ¿Ni dominación ni monotonía?
Fatou sigue acotando un lado, y la igualdad puede fallar de verdad: la
joroba fugitiva $f_n = n\mathbf
1_{\intoo0{1/n}}$ tiene $\int f_n = 1$ pero $f_n \to 0$ \omterm{def:b3:lebesgue:l1}{en casi todo
punto}. La dominación es exactamente lo que impide a la masa escapar al
infinito, vertical u horizontalmente.
\end{method}

\section{Riemann frente a Lebesgue}

\begin{theorem}[Criterio de Lebesgue]
\label{thm:b3:lebesgue:riemann}
Sea $f \colon \intcc ab \to \R$ \emph{acotada}. Entonces $f$ es
\omterm{def:b3:lebesgue:l1}{integrable} Riemann si y solo si $f$ es
\omterm{def:b3:topology:continuity}{continua} en $\lambda$-casi todo
punto; y en tal caso $f$ es \omterm{def:b3:lebesgue:l1}{integrable}
Lebesgue y ambas integrales coinciden.\index{integral de Riemann frente
a la de Lebesgue}
\end{theorem}

\begin{proof}
Para una subdivisión $\sigma = (a = x_0 < \dots < x_N = b)$, sean
$L_\sigma$ y $U_\sigma$ las funciones escalonadas iguales, en cada
$\intoo{x_{i-1}}{x_i}$, a $m_i = \inf_{[x_{i-1}, x_i]}f$ y a
$M_i = \sup$; las sumas de Darboux son sus integrales (Riemann y
Lebesgue coinciden sobre las funciones escalonadas, y ambas dan $\sum
m_i\Delta x_i$). Tómese una sucesión de subdivisiones $\sigma_n$, cada
una refinando la anterior, de paso $\to 0$, cuyas sumas de Darboux
converjan a las integrales de Darboux inferior y superior de $f$. Los
refinamientos hacen $L_{\sigma_n}$ no decreciente y $U_{\sigma_n}$ no
creciente puntualmente fuera del conjunto numerable $D$ de todos los
puntos de división; llámense $\ell$ y $u$ a los límites
(\omterm{def:b3:lebesgue:measurable}{medibles},
la~\cref{prop:b3:lebesgue:stability}). Para $x \notin
D$, escribiendo $I_n(x)$ para el intervalo abierto de $\sigma_n$ que
contiene $x$: $\ell(x) = \sup_n\inf_{I_n(x)}f$ y $u(x) =
\inf_n\sup_{I_n(x)}f$; y como los pasos tienden a $0$, estas son las
\emph{envolventes} inferior y superior de $f$ en $x$ ---
$u(x) - \ell(x)$ es la oscilación de $f$ en $x$ ---, de modo que
\emph{$\ell(x) = u(x)$ si y solo si $f$ es
\omterm{def:b3:topology:continuity}{continua} en $x$}. Por convergencia
monótona o dominada (acotada, intervalo finito):
\[
\int_{\intcc ab}\ell\,\dd\lambda = \lim_n\int L_{\sigma_n}
= \underline{\int}f,
\qquad
\int_{\intcc ab}u\,\dd\lambda = \overline{\int}f .
\]
$f$ \omterm{def:b3:lebesgue:l1}{integrable} Riemann $\iff$
$\underline\int f =
\overline\int f$ $\iff$ $\int(u - \ell) = 0$ $\iff$ $u = \ell$
a.e.\ ($u - \ell \geq 0$; \cref{exo:b3:lebesgue:5}) $\iff$ $f$
\omterm{def:b3:topology:continuity}{continua} \omterm{def:b3:lebesgue:l1}{en casi todo punto}. En tal
caso $\ell \leq f \leq u$ con $\ell =
u$ \omterm{def:b3:lebesgue:l1}{en casi todo punto}: $f$ coincide \omterm{def:b3:lebesgue:l1}{en casi todo punto} con la $\ell$
\omterm{def:b3:lebesgue:measurable}{medible} y, por tanto, es
\omterm{def:b3:lebesgue:measurable}{medible} Lebesgue
(\omterm{def:b3:complete:complete}{completitud} de $\lambda$) con
$\int f\,\dd\lambda = \int\ell\,\dd\lambda = \underline\int f =
\int_a^bf$.
\end{proof}

\begin{example}\label{ex:b3:lebesgue:riemannexamples}
$\mathbf 1_\Q$ no es \omterm{def:b3:topology:continuity}{continua} en
ningún punto: no es \omterm{def:b3:lebesgue:l1}{integrable} Riemann ---
pero es trivial para Lebesgue: $\int\mathbf 1_\Q\,\dd\lambda =
\lambda(\Q) = 0$. La función de Thomae (\,$\frac1q$ en los racionales
$\frac pq$, $0$ en el resto) es
\omterm{def:b3:topology:continuity}{continua} exactamente en los
irracionales: es \omterm{def:b3:lebesgue:l1}{integrable} Riemann con
integral $0$. Y las integrales de Riemann \emph{impropias} son una
noción distinta: $\int_0^\infty\frac{\sin x}x\,\dd x$ converge como
límite de $\int_0^A$ (el problema de fin de semana la calcula: vale
$= \frac\pi2$), pero $\frac{\sin x}x \notin L^1(\intoo0{+\infty})$: la
integral absoluta diverge como la serie armónica
(el~\cref{exo:b3:lebesgue:6}). La teoría de Lebesgue cambia la
convergencia condicional por teoremas de límite robustos.
\end{example}

\section{Integrales con parámetros}

En toda la sección, $(X, \mathcal A, \mu)$ es un
\omterm{def:b3:measure:measure}{espacio de medida}, $T$ un espacio
métrico (el parámetro) y $f \colon T \times X \to \C$ con $f(t, \cdot)$
\omterm{def:b3:lebesgue:l1}{integrable} para cada $t$; póngase
$F(t) = \int_X
f(t, x)\,\dd\mu(x)$.

\begin{theorem}[Continuidad]\label{thm:b3:lebesgue:paramcont}
Supóngase que $t \mapsto f(t,x)$ es
\omterm{def:b3:topology:continuity}{continua} en $t_0$ para casi todo
$x$, y que existe una $g$ \omterm{def:b3:lebesgue:l1}{integrable} con
$\abs{f(t,x)} \leq
g(x)$ para todo $t$ de un \omterm{def:b3:topology:topology}{entorno} de
$t_0$ y casi todo $x$. Entonces $F$ es
\omterm{def:b3:topology:continuity}{continua} en $t_0$.\index{integrales
con parámetro}
\end{theorem}

\begin{proof}
Para toda sucesión $t_n \to t_0$: $f(t_n, \cdot) \to f(t_0,
\cdot)$ \omterm{def:b3:lebesgue:l1}{en casi todo punto}, dominada por $g$: la convergencia dominada
da $F(t_n) \to F(t_0)$; y en los espacios métricos basta la
\omterm{def:b3:topology:continuity}{continuidad} secuencial
(\cref{rem:b3:topology:sequences}).
\end{proof}

\begin{theorem}[Derivación bajo el signo integral]
\label{thm:b3:lebesgue:paramdiff}
Sea $T$ un intervalo abierto de $\R$. Supóngase que, para casi todo $x$,
$t \mapsto f(t,x)$ es derivable en $T$, con
\[
\Bigl|\frac{\partial f}{\partial t}(t, x)\Bigr| \leq g(x)
\quad \text{para todos } t \in T,\ \text{en casi todo punto } x,
\]
$g$ \omterm{def:b3:lebesgue:l1}{integrable}. Entonces $F$ es derivable
en $T$ con $F'(t)
= \int_X \frac{\partial f}{\partial t}(t, x)\,\dd\mu(x)$.
\end{theorem}

\begin{proof}
Fíjense $t$ y $h_n \to 0$: los cocientes incrementales
\[
\varphi_n(x) = \frac{f(t + h_n, x) - f(t, x)}{h_n}
\longrightarrow \frac{\partial f}{\partial t}(t,x)
\quad\text{en casi todo punto},
\]
y la desigualdad del valor medio acota $\abs{\varphi_n(x)} \leq
\sup_{s}\abs{\partial_tf(s,x)} \leq g(x)$: se aplica la convergencia
dominada, y
$\frac{F(t + h_n) - F(t)}{h_n} = \int\varphi_n \to
\int\partial_t f(t, \cdot)$.
\end{proof}

\begin{example}[La función Gamma]\label{ex:b3:lebesgue:gamma}
Para $t > 0$, póngase\index{función Gamma}
\[
\Gamma(t) = \int_0^{+\infty} x^{t-1}\eu^{-x}\,\dd x .
\]
La integral converge: cerca de $0$, $x^{t-1}$ es
\omterm{def:b3:lebesgue:l1}{integrable} ($t >
0$); y en el infinito, $x^{t-1}\eu^{-x} \leq C\eu^{-x/2}$. La
integración por partes (en $[\varepsilon, A]$ y después el límite por
convergencia monótona) da la ecuación funcional $\Gamma(t + 1) =
t\,\Gamma(t)$, de donde $\Gamma(n+1) = n!$: el factorial interpolado. En
todo $\intcc ab \subseteq \intoo0{+\infty}$,
$\partial_t\bigl(x^{t-1}\eu^{-x}\bigr) = \ln
x\cdot x^{t-1}\eu^{-x}$ está dominada por $\abs{\ln x}(x^{a-1} +
x^{b-1})\eu^{-x}$, que es \omterm{def:b3:lebesgue:l1}{integrable}:
$\Gamma$ es $\mathcal C^1$ y, por inducción, $\mathcal C^\infty$, con
$\Gamma^{(k)}(t) =
\int_0^\infty(\ln x)^kx^{t-1}\eu^{-x}\dd x$. El valor
$\Gamma(\frac12) = \sqrt\pi$ es la integral gaussiana disfrazada
(el~\cref{pb:b3:lebesgue:1}).
\end{example}

\begin{omfigure}
\begin{tikzpicture}
\begin{axis}[omaxis, width=10.6cm, height=5.4cm,
    xmin=0, xmax=25, ymin=-0.25, ymax=1.05,
    xtick={0, 3.1416, 6.2832, 9.4248, 12.566, 15.708, 18.850,
      21.991, 25.133},
    xticklabels={$0$, $\pi$, $2\pi$, $3\pi$, $4\pi$, $5\pi$,
      $6\pi$, $7\pi$, $8\pi$},
    ytick={0, 1}]
  \addplot[omDef!40] coordinates {(0,0) (25,0)};
  \addplot[omThm, thick, domain=0.02:25, samples=600]
    {sin(deg(x))/x};
  \addplot[omDef, dashed, domain=0.02:25, samples=200]
    {1/x};
  \addplot[omDef, dashed, domain=0.02:25, samples=200]
    {-1/x};
\end{axis}
\end{tikzpicture}

{\small El integrando $\frac{\sin x}x$: la integral impropia
$\int_0^\infty$ converge por cancelación alternada entre los arcos, pero
las áreas $\abs{\cdot}$ de los arcos se comportan como $\frac2{\pi k}$
--- una serie armónica: $\frac{\sin x}x
\notin L^1$. La \omterm{def:b3:lebesgue:l1}{integrabilidad} de Lebesgue
es \omterm{def:b3:lebesgue:l1}{integrabilidad} absoluta.}
\end{omfigure}

\section{Ejercicios}

\begin{exercise}[$\star$]\label{exo:b3:lebesgue:1}
(a) Demostrar que una función monótona $\R \to \R$ es
\omterm{def:b3:lebesgue:measurable}{medible} Borel, y que una derivada
(de una función derivable en todo punto) es
\omterm{def:b3:lebesgue:measurable}{medible} Borel. (b) Demostrar que
$f \colon X \to \R$ es \omterm{def:b3:lebesgue:measurable}{medible} si y
solo si $\{f > q\}
\in \mathcal A$ para todo $q$ \emph{racional}.
\end{exercise}

\begin{exercise}[$\star$]\label{exo:b3:lebesgue:2}
Calcular, con justificación \omterm{def:b3:complete:complete}{completa}:
\[
\lim_{n\to\infty}\int_0^{+\infty}
\frac{\cos x}{(1 + x/n)^{n}}\,\dd x,
\qquad
\lim_{n\to\infty}\int_0^1 \frac{n\,x^{n-1}}{1 + x}\,\dd x .
\]
\emph{(Para la segunda: sustitúyase $u = x^n$ antes de dominar.)}
\end{exercise}

\begin{exercise}[$\star\star$]\label{exo:b3:lebesgue:3}
(a) Exhibir una desigualdad estricta en el lema de Fatou. (b) Exhibir
$f_n \to 0$ puntualmente con $\int f_n = 1$ de tres maneras: escape en
altura, en anchura y al infinito. ¿Qué única hipótesis de la
convergencia dominada viola cada una? (c) Demostrar que en el lema de
Fatou no se puede sustituir $\liminf$ por $\limsup$ en ninguno de los
dos miembros.
\end{exercise}

\begin{exercise}[$\star\star$]\label{exo:b3:lebesgue:4}
(a) Demostrar que $\displaystyle\int_0^{+\infty}\frac{x}{\eu^x -
1}\,\dd x = \sum_{n\geq1}\frac1{n^2} = \frac{\pi^2}6$
\emph{(desarróllese $\frac1{\eu^x - 1}$ en serie geométrica e intégrese
término a término --- ¿qué teorema lo permite?)}. (b) (El sueño del
estudiante de segundo) Demostrar
$\displaystyle\int_0^1 x^{-x}\,\dd x = \sum_{n\geq1}n^{-n}$.
\emph{(Escríbase $x^{-x} = \eu^{-x\ln x} = \sum_k\frac{(-x\ln
x)^k}{k!}$ y calcúlese $\int_0^1(-x\ln x)^k\dd x$ sustituyendo
$x = \eu^{-u/(k+1)}$ y reconociendo $\Gamma$.)}
\end{exercise}

\begin{exercise}[$\star\star$]\label{exo:b3:lebesgue:5}
(a) Demostrar que $f \geq 0$
\omterm{def:b3:lebesgue:measurable}{medible} con $\int f\,\dd\mu = 0$
cumple $f = 0$ \omterm{def:b3:lebesgue:l1}{en casi todo punto}. \emph{(Considérese $\{f \geq 1/n\}$ y
la desigualdad de Markov: $\mu(\{f \geq a\}) \leq \frac1a\int f$ ---
demuéstrese.)} (b) Demostrar que una $f$
\omterm{def:b3:lebesgue:l1}{integrable} es finita \omterm{def:b3:lebesgue:l1}{en casi todo punto}.
(c) Demostrar que si $\int_A f\,\dd\mu = 0$ para \emph{todo} $A$
\omterm{def:b3:lebesgue:measurable}{medible}, entonces $f = 0$ \omterm{def:b3:lebesgue:l1}{en casi
todo punto}.
\end{exercise}

\begin{exercise}[$\star\star$]\label{exo:b3:lebesgue:6}
(a) Aplíquese el~\cref{thm:b3:lebesgue:riemann} para decidir la
\omterm{def:b3:lebesgue:l1}{integrabilidad} Riemann de: $\mathbf 1_\Q$;
la función de Thomae; $\mathbf
1_K$ para $K$ un conjunto de Cantor gordo (el~\cref{exo:b3:measure:5}).
(b) Demostrar que $\int_1^{+\infty}\abs{\frac{\sin x}x}\,\dd x =
+\infty$, mientras que $\lim_{A\to\infty}\int_1^A\frac{\sin
x}x\,\dd x$ existe \emph{(intégrese por partes)}: convergencia impropia
sin \omterm{def:b3:lebesgue:l1}{integrabilidad}.
\end{exercise}

\begin{exercise}[$\star\star$]\label{exo:b3:lebesgue:7}
Justificar que $F(t) = \int_0^{+\infty}\eu^{-x^2}\cos(tx)\,\dd x$ es
$\mathcal C^1$ en $\R$ y cumple $F'(t) = -\frac t2F(t)$ \emph{(intégrese
por partes)}; deducir $F(t) =
F(0)\,\eu^{-t^2/4}$. (Con $F(0) = \frac{\sqrt\pi}2$ del problema de fin
de semana: la gaussiana es esencialmente su propia transformada de
Fourier --- el~\cref{ch:b3:fouriertransform} sistematizará esto.)
\end{exercise}

\begin{exercise}[$\star\star\star$]\label{exo:b3:lebesgue:8}
(Frullani) Sean $0 < a < b$. Demostrar
\[
\int_0^{+\infty}\frac{\eu^{-ax} - \eu^{-bx}}{x}\,\dd x =
\ln\frac ba,
\]
escribiendo el integrando como $\int_a^b \eu^{-xt}\,\dd t$ y
justificando el intercambio mediante la teoría no negativa
(el~\cref{cor:b3:lebesgue:additivity} en forma
\omterm{def:b3:topology:continuity}{continua} --- anticípese Tonelli, o
córtese $[a,b]$ en $n$ partes iguales y pásese al límite).
\end{exercise}

\begin{exercise}[$\star\star$]\label{exo:b3:lebesgue:9}
Sea $f \geq 0$ \omterm{def:b3:lebesgue:measurable}{medible} sobre
$(X, \mathcal A, \mu)$. Demostrar que $\nu(A) = \int_A f\,\dd\mu$ define
una \omterm{def:b3:measure:measure}{medida} (de
\emph{densidad}\index{densidad de una medida} $f$ respecto de $\mu$), y
que $\int g\,\dd\nu = \int gf\,\dd\mu$ para toda $g \geq 0$
\omterm{def:b3:lebesgue:measurable}{medible} \emph{(demuéstrese para
indicadores, después para \omterm{def:b3:lebesgue:simple}{funciones
simples} y después por convergencia monótona --- la máquina estándar)}.
\end{exercise}

\begin{exercise}[$\star\star\star$]\label{exo:b3:lebesgue:10}
(Un fallo al estilo de Weierstrass) Defínase $f(t) =
\int_0^{+\infty}\frac{\sin(tx)}{x(1 + x^2)}\,\dd x$. (a) Demostrar que
$f$ está bien definida y es
\omterm{def:b3:topology:continuity}{continua} en $\R$, y que
$\mathcal C^1$ con $f'(t) = \int_0^\infty\frac{\cos(tx)}{1 +
x^2}\dd x$ para todo $t$ --- pero que derivar \emph{otra vez} bajo el
signo integral es ilegítimo. (b) Admitiendo que
$f'(t) = \frac\pi2\eu^{-t}$ para $t > 0$ (se demuestra en
el~\cref{ch:b3:residues}), ¿cuánto vale
$\int_0^\infty\frac{x\sin(tx)}{1+x^2}\dd x$ para $t > 0$, y por qué su
fórmula confirma el fallo de (a)?
\end{exercise}

\begin{exercise}[$\star\star$]\label{exo:b3:lebesgue:11}
(Lema de Scheffé) Sean $f_n, f \geq 0$
\omterm{def:b3:lebesgue:l1}{integrables} con $f_n \to f$ \omterm{def:b3:lebesgue:l1}{en casi todo
punto} y $\int f_n \to \int f$. (a) Demostrar que
$\int\abs{f_n - f} \to 0$. \emph{(Aplíquese la convergencia dominada a
$g_n = (f - f_n)^+ \leq f$ y escríbase
$\int\abs{f_n - f} = 2\int g_n - \int(f - f_n)$.)} (b) Demostrar con un
ejemplo que la hipótesis $\int f_n \to \int
f$ no puede suprimirse (una joroba que se desliza o que se concentra), y
que la conclusión falla para $f_n$ con signo sin control en valor
absoluto: $f_n = n\mathbf 1_{\intoc0{1/n}} -
n\mathbf 1_{\intoc{-1/n}0}$ cumple $f_n \to 0$ \omterm{def:b3:lebesgue:l1}{en casi todo punto} y
$\int f_n
= 0 \to 0$, y sin embargo $\int\abs{f_n} = 2$. (c) Aplicación
(\omterm{ex:b3:lebesgue:gamma}{densidades}): si \omterm{ex:b3:lebesgue:gamma}{densidades} de probabilidad cumplen $p_n
\to p$ \omterm{def:b3:lebesgue:l1}{en casi todo punto}, entonces automáticamente
$\int\abs{p_n - p} \to 0$: la convergencia puntual de las \omterm{ex:b3:lebesgue:gamma}{densidades} es
convergencia $L^1$ --- una mejora de convergencia gratis.
\end{exercise}

\begin{exercise}[$\star\star$]\label{exo:b3:lebesgue:12}
Límites clásicos, con justificación \omterm{def:b3:complete:complete}{completa} mediante convergencia
monótona o dominada:
\[
\text{(a)}\ \lim_{n\to\infty}\int_0^n\Bigl(1 -
\frac xn\Bigr)^n\eu^{x/2}\,\dd x, \qquad
\text{(b)}\ \lim_{n\to\infty}\int_0^1\frac{n\,x^{n-1}}{1 +
x}\,\dd x,
\]
\[
\text{(c)}\ \lim_{n\to\infty}\int_0^\infty
\frac{\dd x}{(1 + x/n)^n\,x^{1/n}} .
\]
\emph{(Para (a): $(1 - x/n)^n \nearrow \eu^{-x}$ para $x$ fijo ---
demuéstrese la monotonía mediante $\log$; para (b), intégrese por partes
o sustitúyase $x = u^{1/n}$ e identifíquese una concentración en el
borde; para (c), búsquese un dominador
\omterm{def:b3:lebesgue:l1}{integrable} válido para todo $n \geq 2$
partiendo en $x =
1$.)}
\end{exercise}

\section{Problema: dos integrales célebres}

\begin{problem}[{Problema de fin de semana --- la integral gaussiana y
la integral de Dirichlet, solo con parámetros}]
\label{pb:b3:lebesgue:1}
Dos integrales gobiernan el análisis aplicado:
\[
G = \int_{-\infty}^{+\infty}\eu^{-x^2}\dd x = \sqrt\pi,
\qquad
D = \int_0^{+\infty}\frac{\sin x}{x}\,\dd x = \frac\pi2
\]
(la segunda como integral impropia,
el~\cref{ex:b3:lebesgue:riemannexamples}). Demostraremos ambas usando
únicamente las herramientas de este capítulo.

\medskip
\noindent\textbf{Parte I --- La gaussiana.} Para $t \geq 0$, póngase
\[
A(t) = \Bigl(\int_0^t\eu^{-x^2}\dd x\Bigr)^{2},
\qquad
B(t) = \int_0^1\frac{\eu^{-t^2(1 + x^2)}}{1 + x^2}\,\dd x .
\]
\begin{enumerate}
  \item Justificar que $A$ y $B$ son $\mathcal C^1$ en
  $\intoo0{+\infty}$, calcular $A'$ y $B'$ y demostrar que
  $A'(t) + B'(t) = 0$. \emph{(En $B'$, sustitúyase $u =
        tx$.)} \item Calcular $A(0) + B(0)$ y
  $\lim_{t\to+\infty}(A + B)(t)$ --- justifíquese el paso al límite bajo
  la integral en $B$. \item Concluir $\int_0^\infty \eu^{-x^2}\dd x =
        \frac{\sqrt\pi}2$, de donde $G = \sqrt\pi$, y deducir
  $\Gamma(\tfrac12) = \sqrt\pi$ \emph{(sustitúyase $x =
        u^2$ en $\Gamma(\frac12)$)}.
\end{enumerate}

\medskip
\noindent\textbf{Parte II --- La integral de Dirichlet.} Para $t
\geq 0$, póngase
\[
F(t) = \int_0^{+\infty}\eu^{-tx}\,\frac{\sin x}{x}\,\dd x .
\]
\begin{enumerate}[resume]
  \item Demostrar que la integral que define $F(t)$ converge para todo
  $t > 0$ como integral de Lebesgue, y para $t = 0$ como integral
  impropia; demostrar que $D =
        \lim_{A\to\infty}\int_0^A\frac{\sin x}x\dd x$ existe
  \emph{(intégrese por partes en $[\pi, A]$)}. \item Demostrar que $F$
  es $\mathcal C^1$ en $\intoo0{+\infty}$ con
        \[
        F'(t) = -\int_0^{+\infty}\eu^{-tx}\sin x\,\dd x
        = -\frac{1}{1 + t^2}
        \]
        \emph{(dominación en $[t_0, \infty)$ para cada $t_0 >
        0$; la última integral, por dos integraciones por partes o con
        exponenciales complejas)}. \item Demostrar que $F(t) \to 0$
        cuando $t \to +\infty$ y deducir $F(t)
        = \frac\pi2 - \arctan t$ en $\intoo0{+\infty}$. \item El punto
        delicado: $D = \lim_{t\to0^+}F(t)$. Demuéstrese mediante un
        control uniforme de la cola: para $0 \leq t \leq 1$ y
        $A \geq \pi$, intégrese por partes para ver que
        \[
        \Bigl|\int_A^{+\infty}\eu^{-tx}\frac{\sin
        x}x\,\dd x\Bigr| \leq \frac{C}{A}
        \]
        con $C$ independiente de $t$ \emph{(derívese $\frac{\eu^{-tx}}x$
        y acótese $\abs{\cos}$ por $1$; obsérvese que
        $t\eu^{-tx} \leq 1/x\cdot(tx\eu^{-tx})$ con
        $\sup_{u\geq0}u\eu^{-u} < 1$)}; pártase después $F(t) - D$ en
        $[0, A]$ (donde se aplica la convergencia dominada cuando
        $t \to 0$) y
        $[A, \infty)$.
  \item Concluir: $D = \frac\pi2$.
\end{enumerate}

\medskip
\noindent\textbf{Parte III --- Dividendos.}
\begin{enumerate}[resume]
  \item Calcular $\int_0^{+\infty}\frac{\sin^2x}{x^2}\,\dd x$
  \emph{(intégrese por partes y redúzcase a $D$ mediante $\sin
        2x = 2\sin x\cos x$)}. \item Calcular $\int_0^{+\infty}\frac{1 -
        \cos x}{x^2}\,\dd x$ y comprobar la coherencia de los dos
  resultados. \item Para $a > 0$, calcular
  $\int_0^{+\infty}\frac{\sin(ax)}x\dd x$ y
  $\int_{-\infty}^{+\infty}\eu^{-ax^2}\dd x$, y regístrense las reglas
  de cambio de escala (serán los caballos de batalla de
        \cref{ch:b3:fouriertransform}).
  \item Explicar con precisión por qué $D$ no podía tratarse
  directamente por convergencia dominada en $t = 0$ (no hay dominador
  \omterm{def:b3:lebesgue:l1}{integrable} en $[0,1]\times[0,\infty)$), y
  por qué el corte de la cola de la pregunta 7 es el sustituto honesto
  --- este patrón («\omterm{def:b3:lebesgue:l1}{integrabilidad} uniforme
  de las colas») reaparece a lo largo de todo el análisis.
\end{enumerate}

\medskip
\noindent\textbf{Parte IV --- La \omterm{ex:b3:lebesgue:gamma}{función
Gamma} según Bohr y Mollerup.} La función $\Gamma$
(el~\cref{ex:b3:lebesgue:gamma}) cumple $\Gamma(1) = 1$ y
$\Gamma(x+1) = x\Gamma(x)$ --- pero también lo hacen infinitas funciones
más (multiplíquese por cualquier ondulación $1$-periódica). Una sola
condición de convexidad fija $\Gamma$ de manera única, y sus identidades
más profundas caen entonces mecánicamente. Una función positiva $f$
sobre un intervalo es \emph{logarítmicamente convexa} si $\log f$ es
convexa.
\begin{enumerate}[resume]
  \item Demostrar que la convexidad logarítmica implica la convexidad,
  que los productos de funciones logarítmicamente convexas y sus
  composiciones con aplicaciones afines son logarítmicamente convexos y
  --- mediante la desigualdad de Hölder para dos funciones
  $\int\abs{uv} \leq
        \bigl(\int\abs u^p\bigr)^{1/p}\bigl(\int\abs
        v^q\bigr)^{1/q}$, demostrada directamente a partir de la
  desigualdad de Young --- que $\Gamma$ es logarítmicamente convexa en
        $\intoo0\infty$.
  \item (Lema de las pendientes) Sea $g$ convexa en $\intoo0\infty$ con
  $g(n+1) - g(n) = \log n$ para todo entero $n
        \geq 1$. Para $x \in \intoc01$ y $n \geq 2$, compárense las
  pendientes de $g$ sobre $[n-1, n]$, $[n, n+x]$ y $[n, n+1]$, y
  dedúzcase
        \[
        x\log(n-1) \;\leq\; g(n + x) - g(n) \;\leq\; x\log n .
        \]
  \item (Bohr--Mollerup) Sea $f > 0$ con $f(1) = 1$, $f(x+1) = xf(x)$ y
  $\log f$ convexa. Desenrollando la recursión en
  $f(n + x) = x(x+1)\cdots(x + n -
        1)\,f(x)$ y $f(n) = (n-1)!$, dedúzcase de la pregunta 14 que,
  para $x \in \intoc01$,
        \[
        f(x) = \lim_{n\to\infty}
        \frac{n!\,n^x}{x(x+1)\cdots(x+n)} :
        \]
        $f$ es única, luego $f = \Gamma$, y vale la \emph{fórmula límite
        de Gauss} (extiéndase a todo $x > 0$ mediante la recursión).
        \item Defínase la \emph{función Beta} $B(x, y) =
        \int_0^1t^{x-1}(1-t)^{y-1}\,\dd t$ ($x, y > 0$). Demuéstrense la
        convergencia, la recursión $B(x+1, y) =
        \frac{x}{x+y}\,B(x, y)$ (intégrese por partes) y
        $B(1, y) = \frac1y$.
  \item Demostrar que $x \mapsto B(x, y)$ es logarítmicamente convexa
  (Hölder de nuevo) y aplíquese Bohr--Mollerup a
        \[
        f(x) = \frac{B(x, y)\,\Gamma(x + y)}{\Gamma(y)}
        \]
        para concluir la \emph{fórmula de Euler}: $B(x, y) =
        \dfrac{\Gamma(x)\Gamma(y)}{\Gamma(x+y)}$ --- sin ninguna
        integral doble. \item Calcular $B(\frac12, \frac12)$
        directamente (sustitúyase $t
        = \sin^2\theta$) y deducir $\Gamma(\frac12) =
        \sqrt\pi$: la integral gaussiana de la Parte I, recuperada por
        pura convexidad. Compárense las dos demostraciones en una frase
        cada una. \item (Duplicación de Legendre) Demostrar que
        \[
        g(x) = \frac{2^{x-1}}{\sqrt\pi}\,
        \Gamma\Bigl(\frac x2\Bigr)
        \Gamma\Bigl(\frac{x+1}2\Bigr)
        \]
        cumple las tres hipótesis de Bohr--Mollerup y concluir
        $g = \Gamma$, es decir, $\Gamma(2z) =
        \frac{2^{2z-1}}{\sqrt\pi}\,\Gamma(z)\,\Gamma(z +
        \tfrac12)$ para todo $z > 0$. \item Deducir la forma cerrada
        $\Gamma\bigl(n + \tfrac12\bigr)
        = \dfrac{(2n)!}{4^n\,n!}\sqrt\pi$ y demostrar, mediante el lema
        de las pendientes aplicado a $\log\Gamma$ alrededor de enteros
        grandes, la asintótica
        \[
        \frac{\Gamma(n + \frac12)}{\Gamma(n)\,\sqrt n}
        \longrightarrow 1 .
        \]
  \item Combinar las dos últimas preguntas en la asintótica del
  coeficiente binomial central
        \[
        \binom{2n}{n} \sim \frac{4^n}{\sqrt{\pi n}},
        \]
        y comprobarla numéricamente para $n = 10$ ($\binom{20}{10} =
        184756$, frente a $4^{10}/\sqrt{10\pi} \approx
        187079$: cociente $\approx 0.988$). \item (Síntesis) La
        constante $\sqrt\pi$ ha aparecido ya como integral gaussiana
        (Parte I), como $B(\frac12,
        \frac12)$ (pregunta 18) y dentro de la duplicación (pregunta
        19); y la estimación del binomial central anticipa tanto
        Stirling (el problema de fin de semana del~\cref{ch:b3:product})
        como de Moivre--Laplace. Trácese el mapa de las conexiones: qué
        enunciados equivalen a cuáles, y qué aporta cada técnica ---
        derivación bajo el signo integral frente a convexidad --- que la
        otra no puede aportar.
\end{enumerate}

\medskip
\noindent\textbf{Parte V --- Tres dividendos más.}
\begin{enumerate}[resume]
  \item (Wallis, mediante Beta) Para $p > -1$, sustitúyase $t =
        \sin^2\theta$ para ver que
        \[
        W_p = \int_0^{\pi/2}\sin^p\theta\,\dd\theta
        = \frac12\,B\Bigl(\frac{p+1}2, \frac12\Bigr),
        \]
        y dedúzcase de la recursión de Beta (pregunta 16) que
        $W_{n+2} = \frac{n+1}{n+2}\,W_n$ para enteros $n \geq
        0$. Calcúlense $W_{2n}$ y $W_{2n+1}$ en forma cerrada,
        demuéstrese $W_{2n+1}/W_{2n} \to 1$ por encaje y concluir con el
        \emph{producto de Wallis}
        \[
        \frac\pi2 = \lim_{n\to\infty}
        \prod_{k=1}^{n}\frac{4k^2}{4k^2 - 1} .
        \]
  \item (La gaussiana se encuentra con una frecuencia) Para $b \in \R$,
  póngase
        \[
        \Phi(b) = \int_{-\infty}^{+\infty}
        \eu^{-x^2}\cos(2bx)\,\dd x .
        \]
        Demostrar que $\Phi$ es $\mathcal C^1$ en $\R$, que una
        integración por partes da la ecuación diferencial
        $\Phi'(b) = -2b\,\Phi(b)$, y concluir
        \[
        \Phi(b) = \sqrt\pi\,\eu^{-b^2} :
        \]
        la gaussiana se reproduce a sí misma bajo esta transformación
        --- la única identidad sobre la que funcionará
        el~\cref{ch:b3:fouriertransform}. \item (Integral de Frullani)
        Para $0 < a < b$, demostrar que
        \[
        \int_0^{+\infty}
        \frac{\eu^{-ax} - \eu^{-bx}}{x}\,\dd x
        = \log\frac ba ,
        \]
        derivando respecto del parámetro $a$ (justifíquese la dominación
        en cada $\intco{a_0}{+\infty}$, $a_0 > 0$, e identifíquese la
        constante haciendo $a \to b$). ¿Dónde necesita exactamente el
        integrando su singularidad evitable en $x = 0$?
\end{enumerate}
\end{problem}
